\section{Definicje i podstawy}

Zanim zostaną omówione konkretne metody kompresji wag modeli sztucznej inteligencji, warto uporządkować podstawowe pojęcia związane z reprezentacją danych numerycznych oraz samą kwantyzacją. W literaturze termin ten pojawia się w wielu odmianach, jednak w kontekście modeli uczenia maszynowego oznacza on przede wszystkim redukcję precyzji zapisu parametrów tak, aby zmniejszyć zużycie pamięci, poprawić przepustowość pamięciową i skrócić czas wykonywania obliczeń \cite{jacob2017quantization,dettmers2023case4bit}.

\subsection{Tensory, parametry i warstwy modelu}

Współczesne modele sieci neuronowych opisuje się najczęściej za pomocą tensorów, czyli wielowymiarowych tablic liczb. W praktyce najważniejsze są dwa rodzaje takich danych: aktywacje, które są wyliczane w trakcie działania modelu, oraz parametry, czyli wagi i biasy zapisane po zakończeniu treningu. To właśnie parametry zajmują znaczną część pamięci modelu i stanowią główny obszar zainteresowania przy jego kompresji.

Model składa się z warstw, a każda warstwa realizuje określoną operację matematyczną. W warstwach gęstych, konwolucyjnych i rekurencyjnych dominują mnożenia macierzy i sumowania wartości, dlatego nawet niewielka zmiana sposobu reprezentacji wag może istotnie wpłynąć zarówno na rozmiar modelu, jak i na jego wydajność obliczeniową. W dalszej części pracy pojęcia \textit{waga}, \textit{parametr} i \textit{wartość wagowa} są używane zamiennie, o ile z~kontekstu nie wynika inaczej.

\subsection{Reprezentacja liczb}

W pamięci komputera liczby mogą być zapisywane w różnych formatach. Najczęściej spotykaną reprezentacją w modelach uczenia maszynowego jest format zmiennoprzecinkowy, np. FP32, FP16 lub BF16. Format FP32 zapewnia dużą precyzję, ale zajmuje 32 bity na jedną liczbę. Z kolei formaty o mniejszej liczbie bitów, takie jak FP16, zmniejszają rozmiar danych, lecz kosztem dokładności i zakresu reprezentowanych wartości.

W przypadku kwantyzacji dąży się zazwyczaj do przejścia z reprezentacji zmiennoprzecinkowej do reprezentacji o niższej precyzji, najczęściej całkowitej. Przykładowo, liczby mogą być przechowywane w postaci INT8 albo INT4. Taki zapis jest znacznie bardziej ekonomiczny pamięciowo, lecz wymaga dodatkowego mechanizmu odwzorowującego wartości ciągłe na dyskretne poziomy kwantyzacji. Mechanizm ten opiera się zwykle na współczynniku skali oraz, w niektórych wariantach, na punkcie zerowym.

\subsection{Formaty niskiej precyzji}

W praktyce współczesnych modeli najczęściej spotyka się trzy grupy formatów niskiej precyzji. Pierwszą są klasyczne liczby całkowite, takie jak INT8 i INT4, które stanowią podstawę wielu metod kwantyzacji po treningu \cite{jacob2017quantization,dettmers2023case4bit}.

Drugą grupę stanowią formaty zmiennoprzecinkowe o ograniczonej liczbie bitów, takie jak FP16 i BF16. Wykorzystuje się je tam, gdzie potrzeba kompromisu między dokładnością a rozmiarem danych. W praktyce BF16 jest szczególnie ważny w nowoczesnym treningu i wnioskowaniu modeli, ponieważ zachowuje szeroki zakres wartości przy niższych kosztach pamięciowych.

Trzecią grupę stanowią formaty wyspecjalizowane, takie jak NF4, opisany w pracach o QLoRA, oraz techniki hybrydowe dla 8-bitowego wnioskowania, spopularyzowane przez LLM.int8() \cite{dettmers2023qlora,dettmers2022llmint8}.

\subsection{Uczenie, wnioskowanie i dostrajanie}

Uczenie modelu polega na takim dopasowywaniu jego parametrów, aby minimalizować błąd względem danych treningowych. W praktyce oznacza to wielokrotne przetwarzanie przykładów wejściowych, obliczanie funkcji straty i aktualizację wag. Wnioskowanie jest natomiast etapem użycia już wytrenowanego modelu do generowania odpowiedzi, prognoz lub klasyfikacji.

W pracy pojawia się również pojęcie dostrajania, czyli fine-tuningu. Oznacza ono dalsze uczenie modelu, który został już wcześniej wstępnie wytrenowany, zwykle na mniejszym i bardziej wyspecjalizowanym zbiorze danych. Taki model przejmuje ogólne własności nabyte wcześniej, ale adaptuje się do nowego zadania.

W kontekście metod kompresji ważne jest też pojęcie kalibracji. Jest to wykorzystanie reprezentatywnego, zwykle niewielkiego zbioru danych do oszacowania parametrów potrzebnych do kwantyzacji, takich jak skala, rozkład aktywacji lub znaczenie poszczególnych kanałów. Kalibracja nie jest pełnym treningiem, ale może znacząco wpływać na końcową jakość skompresowanego modelu.

W literaturze kalibracja często pojawia się razem z podejściem \textit{post-training quantization}, gdzie model nie jest ponownie uczony, lecz tylko dostrajany numerycznie na podstawie ograniczonej próbki danych \cite{frantar2023gptq,NIPS1989_6c9882bb}.

\subsection{Rodzaje zadań uczenia maszynowego}

W pracy pojawiają się trzy podstawowe typy zadań: regresja, klasyfikacja binarna i klasyfikacja wieloklasowa. Regresja polega na przewidywaniu wartości liczbowej, na przykład ceny, ryzyka lub innego parametru ciągłego. Klasyfikacja binarna wybiera jedną z dwóch klas, na przykład tak lub nie. Klasyfikacja wieloklasowa przypisuje przykład do jednej z wielu klas, jak w przypadku rozpoznawania słowa ze ścieżki audio albo klasy przynależności obrazu.

Takie rozróżnienie ma znaczenie nie tylko dla doboru metryki oceny, ale również dla samej konstrukcji modelu. Inne funkcje aktywacji, inne warstwy wyjściowe i inne funkcje straty są naturalne dla regresji, a inne dla klasyfikacji.

\subsection{Warstwy i funkcje aktywacji}

Warstwa Dense, zwana też warstwą gęstą, łączy każdy element wejścia z każdym elementem wyjścia za pomocą macierzy wag. Jest to jedna z najprostszych i najczęściej stosowanych warstw w modelach klasyfikacyjnych oraz regresyjnych.

Warstwy Conv1D i Conv2D realizują operację konwolucji, czyli przesuwania filtra po sygnale jednowymiarowym lub dwuwymiarowym (np. obrazie). Są szczególnie przydatne w analizie danych sekwencyjnych, audio i~obrazów, ponieważ potrafią wykrywać lokalne wzorce.

Warstwa GRU należy do rodziny sieci rekurencyjnych i służy do przetwarzania danych sekwencyjnych, w~których istotna jest kolejność elementów. W przeciwieństwie do prostych warstw gęstych uwzględnia ona stan pamięci z poprzednich kroków.

LayerNormalization to warstwa normalizacyjna, która przekształca wartości wewnątrz próbki tak, aby poprawić stabilność uczenia i ułatwić propagację gradientu. Tego typu warstwy często pojawiają się w głębokich architekturach jako sposób na ustabilizowanie przebiegu treningu.

Funkcja aktywacji określa, jak neurony przekształcają otrzymane sygnały. W pracy pojawiają się między innymi GeLU, Sigmoid i Softmax. GeLU jest gładką funkcją aktywacji często stosowaną w nowoczesnych architekturach. Sigmoid odwzorowuje wartości do przedziału od 0 do 1 i dlatego nadaje się do klasyfikacji binarnej. Softmax przekształca wektor wyjściowy w rozkład prawdopodobieństwa po klasach i jest typową funkcją wyjściową w klasyfikacji wieloklasowej.

GeLU zostało opisane jako \textit{Gaussian Error Linear Unit} i pojawia się w licznych nowoczesnych modelach głębokiego uczenia, dlatego w tej pracy traktuję je jako standardową funkcję aktywacji dla głębokich warstw gęstych \cite{hendrycks2023gaussian}.

\subsection{Metryki oceny modeli}

Do porównywania modeli wykorzystuje się metryki, czyli miary jakości predykcji. Dokładność określa, jaki odsetek wszystkich odpowiedzi został przewidziany poprawnie. Precyzja mówi, jaka część przykładów oznaczonych przez model jako pozytywne rzeczywiście taka była. Czułość informuje, jaka część wszystkich rzeczywiście pozytywnych przykładów została wykryta. Miara F1 łączy precyzję i czułość w jedną miarę harmoniczną, dzięki czemu dobrze nadaje się do oceny modeli przy niezbalansowanych danych.

W zadaniach regresyjnych stosuje się inne miary. MAE, czyli średni błąd bezwzględny, pokazuje przeciętną wartość odchylenia predykcji od wartości rzeczywistej. MSE, czyli średni błąd kwadratowy, silniej karze większe odchylenia. RMSE jest pierwiastkiem z MSE i ma tę samą jednostkę co wartość przewidywana, dzięki czemu łatwiej go interpretować.

\subsection{Pojęcia obliczeniowe}

W rozdziale o implementacji pojawiają się także pojęcia związane z realizacją obliczeń. CPU oznacza procesor centralny, GPU jest procesorem graficznym przystosowanym do równoległego wykonywania wielu prostych operacji, a CUDA jest platformą i zestawem narzędzi umożliwiających wykonywanie obliczeń ogólnego przeznaczenia na kartach NVIDIA. Kernel CUDA to pojedyncza funkcja uruchamiana równolegle przez wiele wątków GPU.

CUDA w tej pracy odnosi się do modelu programowania GPU i do zestawu narzędzi opisanych przez NVIDIA, który pozwala kompilować oraz uruchamiać kod wykonywany na kartach graficznych \cite{nvidia2024cuda}.

Pojęcie \textit{batch size} oznacza liczbę przykładów przetwarzanych w jednym kroku obliczeń. W większych \textit{batchach} (partiach) część operacji da się lepiej zrównoleglić, ale rośnie też zapotrzebowanie na pamięć. W przypadku tej pracy istotny jest również rozmiar bloku, czyli liczba wag grupowanych razem podczas kwantyzacji. Ten parametr wpływa bezpośrednio na stosunek jakości modelu do rozmiaru pliku wynikowego.

\subsection{Istota kwantyzacji}

Najprościej rzecz ujmując, kwantyzacja polega na przypisaniu liczb z przestrzeni rzeczywistej do skończonego zbioru wartości reprezentowalnych w wybranym formacie. Jeżeli oznaczymy oryginalną wartość przez $x$, to proces kwantyzacji można opisać funkcją
\begin{equation}
	q = Q(x),
\end{equation}
gdzie $q$ jest wartością skwantyzowaną, a $Q$ to funkcja kwantyzacji. Operacja odwrotna, czyli dekwantyzacja, przybliża wartość oryginalną na podstawie wartości skwantyzowanej:
\begin{equation}
	\hat{x} = D(q),
\end{equation}
gdzie $\hat{x}$ jest przybliżeniem oryginalnej wartości $x$, a $D$ to funkcja dekwantyzacji.

W prostym wariancie symetrycznym przyjmuje się, że skwantyzowana wartość powstaje przez podzielenie liczby przez skalę $s$, zaokrąglenie i ewentualne ograniczenie zakresu do dozwolonych granic. Dla formatu o $b$~bitach zakres ten można zapisać jako
\begin{equation}
	\left[-2^{b-1},\, 2^{b-1}-1\right].
\end{equation}
W praktyce spotyka się także warianty, w których zakres jest przesunięty lub kodowany w sposób niestandardowy, jednak idea pozostaje taka sama: odwzorować wartości rzeczywiste na mniejszą liczbę poziomów dyskretnych.

\subsection{Skala i punkt zerowy}

W kwantyzacji symetrycznej najważniejszym parametrem jest skala. Umożliwia ona przeliczenie wartości z domeny oryginalnej na domenę skwantyzowaną i z powrotem. Jeśli blok wag ma wartości o największej wartości bezwzględnej równej $\max(|x|)$, to klasyczny dobór skali dla $b$ bitów można zapisać jako
\begin{equation}
	s = \frac{\max(|x|)}{2^{b-1}-1}.
\end{equation}

Wariant asymetryczny wykorzystuje dodatkowo punkt zerowy $z_0$, który umożliwia przesunięcie zakresu kwantyzacji względem zera. Jest to przydatne w przypadku danych, których rozkład nie jest symetryczny. Wtedy dekwantyzacja przyjmuje postać
\begin{equation}
	\hat{x} = s(q - z_0).
\end{equation}

W modelach sieci neuronowych częściej stosuje się jednak kwantyzację symetryczną, ponieważ wagi po treningu zwykle mają rozkład zbliżony do symetrycznego wokół zera. Upraszcza to implementację \cite{frantar2023gptq}, ale potrafi wpłynąć negatywnie na rozmiar pliku i/lub zakres kwantyzacji ze względu na przechowywanie liczb ujemnych.

\subsection{Kwantyzacja per-tensor, per-channel i blokowa}

Sposób wyznaczania skali ma duże znaczenie dla jakości końcowego modelu. W kwantyzacji \textit{per-tensor} dla całego tensora stosuje się jeden wspólny współczynnik skali. Jest to rozwiązanie najprostsze i najtańsze obliczeniowo, ale może być zbyt mało precyzyjne dla dużych warstw o zróżnicowanych wartościach wag.

Dokładniejszym podejściem jest kwantyzacja \textit{per-channel}, w której osobną skalę wyznacza się dla każdego kanału, wiersza lub kolumny tensora, zależnie od implementacji. Taki wariant lepiej dopasowuje się do lokalnych różnic w rozkładzie wag, ale wymaga przechowywania większej liczby współczynników pomocniczych.

W praktycznych implementacjach bardzo często stosuje się także kwantyzację blokową, zwaną również grupową. W tym podejściu tensor dzieli się na mniejsze fragmenty o ustalonym rozmiarze, a dla każdego bloku wyznacza się osobną skalę. Pozwala to osiągnąć kompromis pomiędzy jakością rekonstrukcji a rozmiarem metadanych. Im mniejszy blok, tym dokładniejsze odwzorowanie wartości, ale tym większy narzut pamięciowy. Zależność ta będzie widoczna również w dalszych rozdziałach pracy, gdzie różne rozmiary bloków wpływają na końcowy rozmiar pliku oraz na jakość modelu po dekwantyzacji.

\subsection{Kwantyzacja po treningu i kwantyzacja podczas treningu}

Ze względu na moment wprowadzenia ograniczeń precyzji rozróżnia się dwa podstawowe podejścia. Kwantyzacja po treningu, czyli PTQ (\textit{Post-Training Quantization}), polega na skompresowaniu już wytrenowanego modelu bez ponownego uczenia wszystkich parametrów. Jest to rozwiązanie wygodne, ponieważ nie wymaga kosztownego procesu dostrajania modelu od początku, a przy odpowiednim doborze metody może zachować bardzo dobrą jakość \cite{frantar2023gptq}.

Drugi wariant to kwantyzacja podczas treningu, czyli QAT (\textit{Quantization-Aware Training}). W tym przypadku model jest uczony z uwzględnieniem ograniczeń wynikających z późniejszej niskiej precyzji zapisu. Pozwala to lepiej przygotować sieć na obecność błędów kwantyzacji, ale wymaga znacznie większych nakładów obliczeniowych i bardziej złożonego procesu trenowania.

W niniejszej pracy główny nacisk położony jest na podejście PTQ, ponieważ pozwala ono badać wpływ kompresji bez konieczności ponownego trenowania modeli. Jest to także podejście bliższe praktycznym narzędziom do kompresji modeli, stosowanym po zakończeniu ich uczenia.

\subsection{Błąd kwantyzacji}

Każde uproszczenie reprezentacji numerycznej powoduje błąd kwantyzacji. Można go rozumieć jako różnicę pomiędzy oryginalną wartością $x$ a jej odpowiednikiem po dekwantyzacji $\hat{x}$. Błąd ten można zapisać jako
\begin{equation}
	e = x - \hat{x}.
\end{equation}

W idealnym przypadku błędy te są niewielkie i rozkładają się w sposób, który nie wpływa znacząco na końcowy wynik modelu. W praktyce sytuacja zależy jednak od rozkładu wag, rodzaju warstwy oraz od tego, czy w danych występują wartości odstające. Właśnie dlatego niektóre modele zachowują bardzo dobrą jakość nawet po silnej kompresji, podczas gdy inne ulegają wyraźnej degradacji już przy umiarkowanym skróceniu precyzji.

Z punktu widzenia dalszych rozdziałów istotne jest więc przede wszystkim to, że skuteczna kwantyzacja nie polega wyłącznie na „ucięciu bitów”, lecz na takim dobraniu skali, sposobu grupowania i reprezentacji pomocniczej, aby zminimalizować wpływ błędu rekonstrukcji na zachowanie modelu.

\subsection{Notacja używana dalej w pracy}

W tym i dalszych rozdziałach pojawia się kilka grup symboli, które warto uporządkować już na tym etapie. Najważniejsze z nich związane są z kwantyzacją wag i z zapisem warstw modelu.

Oryginalną wartość liczbową oznaczam jako $x$, a jej wersję po kwantyzacji jako $q$. Wartość odtworzona w~procesie dekwantyzacji ma postać $\hat{x}$. Błąd kwantyzacji oznaczam symbolem $e$. Funkcję kwantyzacji zapisuję jako $Q$, a funkcję dekwantyzacji jako $D$.

W przypadku wag modeli przyjmuję zapis $w$ dla pojedynczej wagi oraz $W$ dla całej macierzy wag. Skwantyzowany odpowiednik macierzy wag oznaczam jako $\hat{W}$, natomiast macierz wejść jako $X$. W rozdziale o GPTQ pojawia się również $H_F^{-1}$, czyli odwrotna macierz Hesjana, wykorzystywana przy analizie wpływu błędu na wyjście warstwy. Gdy w tekście mowa o różnicy pomiędzy wartością oryginalną i przybliżoną w bardziej ogólnym sensie, używany jest symbol $\Delta$.

Przy opisie skali i przesunięcia stosuję symbole $s$, $z$ oraz $z_0$. W zależności od kontekstu $s$ oznacza współczynnik skali używany do odwzorowania liczb z oryginalnej domeny na domenę skwantyzowaną i z powrotem. Symbol $z$~oznacza przesunięcie lub punkt zerowy w zapisie ogólnym, natomiast $z_0$ jest punktem zerowym w zapisie bardziej precyzyjnym, szczególnie wtedy, gdy trzeba wyróżnić wariant asymetryczny.

W częściach dotyczących blokowej kompresji wag pojawiają się także symbole $b$ i $B$. Litera $b$ oznacza liczbę bitów przeznaczonych na zapis jednej wartości, a $B$ rozmiar bloku, czyli liczbę wag grupowanych wspólnie przy wyznaczaniu skali i innych parametrów pomocniczych. Jeżeli tekst odnosi się do granic reprezentacji dla danego formatu, pojawia się również zapis zakresu dopuszczalnych wartości, np. dla INT4 lub INT8.

W opisie warstw i modeli często występują też zapisy związane z architekturą sieci. Dense, Conv1D, Conv2D, GRU oraz LayerNormalization oznaczają konkretne typy warstw, a GeLU, Sigmoid i Softmax oznaczają funkcje aktywacji używane w ich wyjściach. Gdy w tekście mowa o jakości modelu, stosowane są nazwy metryk: dokładność, precyzja, czułość, miara F1, MAE, MSE i RMSE.

W częściach dotyczących implementacji i testów pojawiają się ponadto skróty obliczeniowe i organizacyjne. CPU oznacza procesor centralny, GPU procesor graficzny, CUDA platformę obliczeń równoległych, a kernel CUDA pojedynczą funkcję wykonywaną równolegle na GPU. Pojęcie \textit{batch size} oznacza rozmiar paczki danych przetwarzanych w jednym kroku, a \textit{block size} rozmiar bloku wag wykorzystywanego w kwantyzacji.
